\documentclass[a4paper,10pt]{article}

% ---- Packages ----
\usepackage[utf8]{inputenc}
\usepackage[T1]{fontenc}
\usepackage[default]{lato}
\usepackage{lmodern}
\usepackage{xcolor}
\usepackage{geometry}
\usepackage{tabularx}
\usepackage{titlesec}
\usepackage{enumitem}
\usepackage[hidelinks]{hyperref}
\usepackage{fancyhdr}
\usepackage{fontawesome5}
\usepackage[most]{tcolorbox}
\usepackage{graphicx}

\geometry{left=1.4cm, top=0.8cm, right=1.2cm, bottom=1cm}
\pagestyle{fancy}
\fancyhf{}
\renewcommand{\headrulewidth}{0pt}
\renewcommand{\footrulewidth}{0pt}
\urlstyle{same}
\setlength{\tabcolsep}{0pt}
\DeclareUnicodeCharacter{202F}{~}

% ---- Colonnes pour tabularx ----
\newcolumntype{L}{>{\raggedright\arraybackslash}p{0.72\linewidth}}
\newcolumntype{R}{>{\raggedleft\arraybackslash}p{0.28\linewidth}}

% ---- Formatage des titres ----
\titleformat{\section}{\scshape\large\raggedright}{}{0em}{}[\titlerule]

% ---- Commandes personnalisées & paramètres de liste ----
\newcommand{\resumeSubheading}[4]{%
  \item
    \begin{tabularx}{\linewidth}{L R}
      \textbf{#1} & \textit{\footnotesize #2} \\
      \textit{\footnotesize #3} & \footnotesize #4 \\
    \end{tabularx}
}

\newcommand{\resumeSubHeadingListStart}{%
  \begin{itemize}[leftmargin=0cm, label={}, itemsep=1pt, topsep=0pt]}
\newcommand{\resumeSubHeadingListEnd}{\end{itemize}}

\newcommand{\resumeItemListStart}{\begin{itemize}[leftmargin=1.5em, label=--, noitemsep, topsep=0pt]}
\newcommand{\resumeItemListEnd}{\end{itemize}}
\newcommand{\resumeItem}[1]{\item #1}

\newcommand{\cvsection}[1]{%
  \vspace{3mm}
  \begin{tcolorbox}[colback=gray!20, colframe=white, boxsep=4pt]
    \textbf{\large #1}
  \end{tcolorbox}
  \vspace{2mm}
}

% ---- Informations personnelles ----
\newcommand{\name}{Sami Hanine}
\newcommand{\phone}{+33 7 59 54 09 57}
\newcommand{\emaila}{sami.hanine.1@ens.etsmtl.ca}
\newcommand{\website}{https://samihanine.fr}
\newcommand{\github}{https://github.com/samihanine}
\newcommand{\linkedin}{http://linkedin.com/in/sami-hanine-a685201b9/}

% ---- Document ----
\begin{document}

% En-tête
\noindent
\begin{minipage}[c]{2.6cm}
  {\setlength{\fboxsep}{0pt}\setlength{\fboxrule}{0.8pt}%
   \fcolorbox{gray!60}{white}{\includegraphics[width=2.6cm]{photo-sami.jpg}}}
\end{minipage}%
\hfill
\begin{minipage}[c]{\dimexpr\linewidth-3.0cm\relax}
  \begin{tabularx}{\linewidth}{@{}>{\raggedright\arraybackslash}X >{\raggedleft\arraybackslash}p{6.2cm}@{}}
    \textbf{\Large \name} & \phone \ \faPhone \\
    Ingénieur spécialisé en IA et analyse de données   & \href{mailto:\emaila}{\emaila} \ \faEnvelope \\
    \textit{Diplômé de l'ÉTS et de HEC Montréal} & \href{\website}{samihanine.fr} \ \faGlobe \\
    \textit{Bilingue anglais \& français} & \href{\github}{github.com/samihanine} \ \faGithub \\
    \textit{Paris, France} &  \href{\linkedin}{linkedin.com/in/sami-hanine} \ \faLinkedin
  \end{tabularx}
\end{minipage}

% ---- Formation ----
\cvsection{Formation}
\resumeSubHeadingListStart
  \resumeSubheading
    {Diplôme d'ingénieur -- ÉTS Montréal}{Sept. 2022 -- Fév. 2027}{Spécialité informatique}{}

  \resumeSubheading
    {\hspace{1em}\textbullet\ Membre du club finance}{Sept. 2025 -- Présent}{\hspace{2.5em}Plateforme de trading sur les marchés}{}

  \resumeSubheading
    {\hspace{1em}\textbullet\ Semestre d'échange à Taïwan}{Fév. 2025 -- Juin 2025}{\hspace{2.5em}Axé sur l'apprentissage automatique}{}

  \resumeSubheading
    {Programme entrepreneurial -- HEC Montréal}{Jan. 2024 -- Avr. 2024}{Construction d'une startup en incubateur}{}

  \resumeSubheading
    {Diplôme en informatique -- Université Claude Bernard}{Sept. 2020 -- Juin 2022}{Diplôme universitaire de technologie en informatique}{}
\resumeSubHeadingListEnd

% ---- Expériences ----
\cvsection{Expériences}
\resumeSubHeadingListStart

  \resumeSubheading
    {Stage en gestion de projet d'analyse de données pour la supply chain de Dior}{Sept. 2026 -- Fév. 2027}{Maison Dior - LVMH -- Paris, France}{}
    \resumeItemListStart
    \resumeItem{Recueil des besoins auprès des équipes business et coordination avec l'équipe IT pour leur réalisation.}
    \resumeItem{Rôle de Product Manager (PM) : priorisation, coordination, specifications, documentations, tests et support.}
    \resumeItem{Intégration d'outils IA pour automatiser les processus de la supply chain et faire gagner du temps aux équipes.}
    \resumeItemListEnd
    \vspace{2mm}

  \resumeSubheading
    {Stage dans un laboratoire de recherche en Intelligence Artificielle}{Sept. 2025 -- Déc. 2025}{Croesus Fintech -- Montréal, Canada}{}
    \resumeItemListStart
      \resumeItem{Développement de pipelines ML/LLM pour l'analyse de portefeuilles.}
      \resumeItem{Entraînement de modèles avec \textbf{Unsloth} et \textbf{PyTorch}, orchestrés via \textbf{Databricks} et \textbf{MLflow}.}
      \resumeItem{Amélioration de la précision du modèle de \textbf{17\%} et déploiement accessible depuis le cloud.}
    \resumeItemListEnd
    \vspace{2mm}

  \resumeSubheading
    {Co-fondateur et directeur d'une agence de services informatiques}{Mars 2023 -- Aug. 2026}{Otopio -- Montréal, Canada}{}
    \resumeItemListStart
      \resumeItem{Traduction des besoins clients en solutions techniques évolutives intégrant l'\textbf{IA}.}
      \resumeItem{Développement de produits en \textbf{TypeScript}, hébergés sur \textbf{AWS}, avec des tableaux de bord \textbf{Power BI}.}
      \resumeItem{Collaboration avec \textbf{20+ clients} et maintien d'applications utilisées par des milliers d'utilisateurs.}
    \resumeItemListEnd
    \vspace{2mm}

  \resumeSubheading
    {Stage en analyse de données dans une startup fintech en Californie}{Sept. 2023 -- Déc. 2023}{Finance It Forward -- Los Angeles, États-Unis}{}
    \resumeItemListStart
      \resumeItem{Automatisation de l'analyse de données financières pour la comparaison de prêts.}
      \resumeItem{Extraction et classification de données avec le \textbf{NLP} via \textbf{Hugging Face}, et mise en place d'un \textbf{RAG} avec \textbf{LangChain}.}
      \resumeItem{Analyses entièrement automatisées qui nécessitaient jusqu'à \textbf{8 heures par semaine} aux employés.}
    \resumeItemListEnd
    \vspace{2mm}
\resumeSubHeadingListEnd

% ---- Compétences techniques ----
\cvsection{Compétences techniques}
\begin{itemize}[leftmargin=0.2cm, label={}, itemsep=1pt, topsep=0pt]
  \item \textbf{Intelligence artificielle} : Fine-tuning, Unsloth, NLP, PyTorch, TensorFlow, LLM, Hugging Face, scikit-learn, RAG, bases de données vectorielles, LangChain, LangGraph, agents IA, OpenRouter, MLflow, MLOps
  \item \textbf{Données} : Pandas, NumPy, SQL, prétraitement des données, ETL \& pipelines de données, web scraping, BigQuery, Snowflake, Databricks, Jupyter Notebook, Power BI
  \item \textbf{Code \& Automatisation} : Python, TypeScript, JavaScript, Go, SQL, Java, C\#, C, C++, PHP, VBA, React, Docker, Git, PostgreSQL, MySQL, Firebase, AWS, Linux/Unix
\end{itemize}

\newpage

% ---- Commandes d'espacement page 2 ----
\newcommand{\projListStart}{%
  \begin{itemize}[leftmargin=0cm, label={}, itemsep=10pt, topsep=0pt]}
\newcommand{\projListEnd}{\end{itemize}}

\newcommand{\projItemStart}{\begin{itemize}[leftmargin=1.5em, label=--, itemsep=4pt, topsep=2pt]}
\newcommand{\projItemEnd}{\end{itemize}}

\cvsection{Projets personnels}
\projListStart

  % ---- medic.pub ----
  \resumeSubheading
    {Medic.pub : Logiciel de gestion de cabinet médical}{}{\href{https://www.medic.pub/}{medic.pub}}{2026}
    \projItemStart
      \resumeItem{\textbf{Contexte :} Mission réalisée en tant que prestataire pour Otopio.}
      \resumeItem{\textbf{Objectif :} Permettre la connexion à des machines médicales pour un logiciel de gestion de cabinet.}
      \resumeItem{\textbf{Comment :} Mise en place d'un flux \textbf{ETL} pour collecter les données brutes des machines, les nettoyer, les transformer et les envoyer vers le cloud.}
      \resumeItem{\textbf{Résultat :} Intégration de deux sources de données (\textbf{HL7} et \textbf{EDFv2}) normalisées au format \textbf{FHIR}.}
    \projItemEnd

  % ---- AlgoETS ----
  \resumeSubheading
    {AlgoETS : Automatisation du trading avec l'IA}{}{\href{https://algoets.etsmtl.ca}{algoets.etsmtl.ca}}{2025}
    \projItemStart
      \resumeItem{\textbf{Contexte :} Travail pour le club finance de mon école.}
      \resumeItem{\textbf{Objectif :} Permettre aux utilisateurs de générer des modèles d'automatisation de trading adaptés à leur propre portefeuille.}
      \resumeItem{\textbf{Comment :} Combinaison de modèles de trading statistiques avec un \textbf{LLM} qui convertit des requêtes en langage naturel en paramètres de modèle.}
      \resumeItem{\textbf{Résultat :} Démo fonctionnelle et intégration de fonctionnalités pour une plateforme avec \textbf{100+ utilisateurs}.}
    \projItemEnd

  % ---- Oragius ----
  \resumeSubheading
    {Oragius : Prédiction du niveau d'eau}{}{\href{https://b2bhint.com/fr/company/ca-qc/les-solutions-oragius-inc--1180974140}{oragius.ca}}{2025}
    \projItemStart
      \resumeItem{\textbf{Contexte :} Projet de recherche avec un professeur de l'ÉTS, entrepreneur en résidence au Centech.}
      \resumeItem{\textbf{Objectif :} Prédire les niveaux d'eau à partir d'un ensemble de données statistiques.}
      \resumeItem{\textbf{Comment :} Prévision de séries temporelles avec un \textbf{LSTM} (\textbf{TensorFlow}, Python) sur des données historiques.}
      \resumeItem{\textbf{Résultat :} Atteinte d'une précision de \textbf{59\%} entre les jeux de données d'entraînement et de test.}
    \projItemEnd

  % ---- MTA Region ----
  \resumeSubheading
    {MTARegion : Analyse du comportement utilisateur}{}{\href{https://mtaregion.com/}{mtaregion.com}}{2024}
    \projItemStart
      \resumeItem{\textbf{Contexte :} Mission en freelance pour une entreprise de tourisme offrant des réductions aux utilisateurs dans les commerces locaux.}
      \resumeItem{\textbf{Objectif :} Analyser les données de comportement utilisateur et évaluer la rentabilité.}
      \resumeItem{\textbf{Comment :} Jointure des données financières (\textbf{Stripe}) avec les données utilisateurs dans \textbf{PostgreSQL} pour alimenter des tableaux de bord.}
      \resumeItem{\textbf{Résultat :} Analyse et catégorisation des données de \textbf{10 000+ clients}.}
    \projItemEnd

  % ---- Coordia ----
  \resumeSubheading
    {Coordia : Collecte de données patients}{}{\href{https://www.coordia.ca/}{coordia.ca}}{2024}
    \projItemStart
      \resumeItem{\textbf{Contexte :} Mission réalisée en tant que prestataire pour Otopio.}
      \resumeItem{\textbf{Objectif :} Collecter des données patients pour un logiciel de gestion de cabinet ortho-médical.}
      \resumeItem{\textbf{Comment :} Mise en place de l'analytique \textbf{PostHog} et développement de tableaux de bord \textbf{BI} personnalisés adaptés au produit.}
      \resumeItem{\textbf{Résultat :} \textbf{50 000+} points de données traités, analysés et catégorisés chaque mois.}
    \projItemEnd

  % ---- Nautilus OS ----
  \resumeSubheading
    {Nautilus OS : Système d'exploitation simulé}{}{\href{https://github.com/samihanine/nautilius-os-batch}{github.com/samihanine/nautilius-os-batch}}{2016}
    \projItemStart
      \resumeItem{\textbf{Contexte :} Projet réalisé sur plus d'un an, entre 13 et 14 ans.}
      \resumeItem{\textbf{Objectif :} Simuler un système d'exploitation dans la console.}
      \resumeItem{\textbf{Comment :} Développement d'un système de mémoire en \textbf{C} et de scripts d'automatisation en \textbf{VBA} et \textbf{Batch} pour contrôler l'ordinateur.}
      \resumeItem{\textbf{Résultat :} Création de \textbf{8 applications complètes} dont un explorateur de fichiers, un éditeur de texte et l'accès/contrôle des périphériques.}
    \projItemEnd

\projListEnd

\end{document}